% =====================================================================
% Kapitel 6: ETL-Pipeline
% =====================================================================

\chapter{ETL-Pipeline}

\label{ch:etl}

Die \emph{ETL-Pipeline} -- Extract, Transform, Load -- ist das
\emph{Rückgrat} des Projekts. Sie ist verantwortlich dafür, dass
die Daten der Vendor-Adapter tatsächlich in die Datenbank gelangen
und für Analysen verfügbar werden.

\section{Überblick}

\begin{definition}
\textbf{ETL}
\emph{ETL} steht für:
  \begin{itemize}
    \item \textbf{Extract} -- Daten aus den Quellen (FRED, ECB, Yahoo)
          abrufen.
    \item \textbf{Transform} -- Daten in das richtige Format und Schema
          bringen.
    \item \textbf{Load} -- Daten in die Datenbank schreiben.
  \end{itemize}
\end{definition}

Abbildung~\ref{fig:etl-flow} zeigt den ETL-Flow im Detail.

\begin{figure}[ht]
\centering
\begin{tikzpicture}[
  node distance=0.8cm and 1.2cm,
  box/.style={rectangle, draw, thick, rounded corners, minimum width=3cm, minimum height=0.8cm, align=center, font=\small},
  arrow/.style={-Latex, thick},
]
  \node[box, fill=red!20] (extract) {1. Extract\\\scriptsize Vendor-Adapter\\\scriptsize (FRED/ECB/Yahoo)};
  \node[box, fill=yellow!20, right=of extract] (validate) {2. Validate Raw\\\scriptsize Pandera Schemas};
  \node[box, fill=orange!20, right=of validate] (transform) {3. Transform\\\scriptsize Polars Mapping\\\scriptsize to canonical schema};
  \node[box, fill=cyan!20, right=of transform] (validate2) {4. Validate Clean\\\scriptsize PK, OHLC, ranges};
  \node[box, fill=green!20, below=of validate] (upsert) {5. Upsert\\\scriptsize SQLAlchemy merge()};
  \node[box, fill=blue!20, right=of upsert] (db) {6. Persist\\\scriptsize PostgreSQL};

  \draw[arrow] (extract) -- (validate);
  \draw[arrow] (validate) -- (transform);
  \draw[arrow] (transform) -- (validate2);
  \draw[arrow] (validate2) -- (upsert);
  \draw[arrow] (upsert) -- (db);
\end{tikzpicture}
\caption{Der ETL-Flow: Sechs Schritte, jeweils mit klarer Verantwortung.
  Im aktuellen Projekt sind die Validate-Schritte vorbereitet, aber
  noch nicht vollständig implementiert (geplant für eine DQ-Iteration).}
\label{fig:etl-flow}
\end{figure}

\section{Der Seeder-Pattern}

Listing~\ref{lst:seeder} zeigt das Standard-Pattern zum
\emph{Seed} (Initial-Befüllung) einer Tabelle. Es wird für
\code{vendors}, \code{instruments} und \code{prices\_daily}
verwendet.

\begin{lstlisting}[language=python, caption={Seeder-Pattern: lookup-or-create},label=lst:seeder]
def upsert_instrument(asset_class, name, currency, symbol):
    """Idempotent: creates the instrument if not in DB, else returns id."""
    with session_scope() as s:
        existing = (
            s.query(Instrument)
            .join(InstrumentIdentifier,
                  InstrumentIdentifier.instrument_id == Instrument.instrument_id)
            .filter(
                InstrumentIdentifier.scheme == "YAHOO",
                InstrumentIdentifier.value == symbol,
            )
            .one_or_none()
        )
        if existing is None:
            inst = Instrument(asset_class=asset_class, name=name, currency=currency)
            s.add(inst)
            s.flush()  # gets us inst.instrument_id
            s.add(InstrumentIdentifier(
                instrument_id=inst.instrument_id,
                scheme="YAHOO",
                value=symbol,
            ))
            s.flush()
            return inst.instrument_id
        return existing.instrument_id
\end{lstlisting}

Die Schlüsseleigenschaft ist \emph{Idempotenz}: Wenn der Seeder
zweimal läuft, werden keine Duplikate angelegt. \code{merge()} allein
reicht dafür nicht aus -- wenn das Objekt keine \code{primary key}
hat, versucht SQLAlchemy ein \code{INSERT} statt eines \code{UPDATE}.
Die explizite \code{SELECT}-Suche im Voraus ist robuster.

\section{Das \code{query\_db.py}-Skript}

Das \code{query\_db.py}-Skript ist die zentrale CLI für Daten-Operationen.
Es kann:

\begin{itemize}
  \item die Datenbank mit Demo-Daten füllen (falls leer),
  \item echte Daten von FRED, ECB und Yahoo holen,
  \item die Daten in tabellarischer Form im Terminal anzeigen.
\end{itemize}

Listing~\ref{lst:query-db} zeigt einen vereinfachten Auszug.

\begin{lstlisting}[language=python, caption={Auszug aus \code{scripts/query\_db.py}},label=lst:query-db]
def main() -> int:
    settings = get_settings()
    configure_logging(settings.log_level)

    with session_scope() as s:
        n_obs = s.execute(select(func.count()).select_from(MacroObservation)).scalar() or 0

    if n_obs == 0 and not args.no_seed:
        # No data: try ECB + FRED in sequence
        if settings.fred_api_key:
            print("seeding real FRED data...")
            _seed_from_fred(settings.fred_api_key)
        print("seeding real ECB data (no API key required)...")
        _seed_from_ecb()
        print("seeding real Yahoo Finance data (AAPL, MSFT, ^GSPC)...")
        _seed_from_yahoo()

    # Display all series with their latest values
    series_list = _fetch_series(args.series)
    _print_summary(series_list, tail=args.tail)
    return 0
\end{lstlisting}

Beachten Sie das \emph{Reihenfolge}-Pattern:

\begin{enumerate}
  \item Prüfe, ob die DB bereits Daten hat (\code{n\_obs > 0}).
  \item Falls leer: Versuche ECB (kein Key) und FRED (falls Key
        gesetzt). Fällt auf Demo-Daten zurück, wenn beide fehlschlagen.
  \item Lese alle Serien und zeige sie formatiert an.
\end{enumerate}

\section{Bulk-Seeding mit \code{seed\_paper\_data.py}}

Für die AEGIS-Paper-Replikation (Kapitel~\ref{ch:analytics}) haben
wir ein eigenes Skript gebaut, das 36~Symbole (31~Aktien + 5~Benchmarks)
in einem einzigen Aufruf von \code{yfinance} lädt.

\begin{lstlisting}[language=shell, caption={Bulk-Seed-Skript aufrufen}]
# Einmaliger Aufruf
.venv/bin/python scripts/seed_paper_data.py

# Mit benutzerdefiniertem Zeitraum
.venv/bin/python scripts/seed_paper_data.py --start 2006-01-01

# Ohne Benchmarks
.venv/bin/python scripts/seed_paper_data.py --no-benchmarks
\end{lstlisting}

Das Skript nutzt die \code{fetch\_many}-Bulk-Methode des
Yahoo-Vendors (Listing~\ref{lst:yahoo-fetch}) -- 36~Symbole
werden in einem einzigen \code{yf.download}-Aufruf geholt, was
~12~Sekunden statt 36~einzelner Calls benötigt. Insgesamt
werden 185.220~Zeilen in die \code{prices\_daily}-Tabelle
geschrieben.

\section{DQ-Checks (Vorbereitung)}

Das Projekt hat Skelette für Data-Quality-Checks, die in einer
späteren Iteration implementiert werden sollen. Die
DQ-Tabellen (\code{dq\_runs}, \code{dq\_findings}) existieren bereits.

Folgende Checks sind geplant:

\begin{itemize}
  \item \textbf{PK-Eindeutigkeit}: \code{(instrument\_id, date)}
        darf nicht doppelt vorkommen.
  \item \textbf{OHLC-Konsistenz}: \code{low} $\leq$ \code{open}, \code{close}
        $\leq$ \code{high} (keine \emph{inverted candles}).
  \item \textbf{Keine negativen Preise/Volumes}.
  \item \textbf{Adjusted-Close-Konsistenz}: Wenn \code{corporate\_actions}
        verfügbar ist, muss \code{adjusted\_close} das Produkt aus
        \code{close} und allen Splits/Dividenden sein.
  \item \textbf{Lücken-Detection}: Mehr als 5~aufeinanderfolgende
        Handelstage ohne Eintrag -- ein \emph{DQ-Finding}.
  \item \textbf{Ausreißer-Detection}: Tagesreturn $> 6\sigma$
        gegenüber rollendem 252-Tage-Fenster -- \emph{flag}, nicht
        löschen.
  \item \textbf{Currency-Coverage}: Jedes \code{instrument} hat
        \code{currency} gesetzt.
  \item \textbf{Vendor-Coverage}: Jede Zeile in \code{prices\_daily}
        hat \code{vendor\_id} gesetzt.
\end{itemize}

\section{Idempotenz}

\begin{merke}
Alle Seeder sind \emph{idempotent}. Sie können beliebig oft
ausgeführt werden, ohne Duplikate zu erzeugen. Dies ist entscheidend
für die Inkrementelle-Aktualisierung: ein täglicher Cron-Job ruft
den Seeder auf, und nur neue Daten werden tatsächlich geschrieben.
\end{merke}

Listing~\ref{lst:idempotency} zeigt das zentrale Idempotenz-Pattern.

\begin{lstlisting}[language=python, caption={Idempotenz-Pattern: SELECT vor INSERT},label=lst:idempotency]
def upsert_prices(instrument_id, vendor_id, df) -> int:
    if df.is_empty():
        return 0
    rows = 0
    with session_scope() as s:
        for r in df.iter_rows(named=True):
            s.merge(PricesDaily(
                instrument_id=instrument_id,
                date=r["date"],
                open=r.get("open"),
                high=r.get("high"),
                low=r.get("low"),
                close=r["close"],
                adjusted_close=r.get("adjusted_close"),
                volume=r.get("volume"),
                vendor_id=vendor_id,
            ))
            rows += 1
    return rows
\end{lstlisting}

\section{Failure-Isolation}

\begin{definition}
\textbf{Failure-Isolation}
\emph{Failure-Isolation} bedeutet, dass der Fehler eines Vendors
\emph{nicht} den gesamten Pipeline-Lauf abbricht. Wenn Yahoo
gerade ein Problem hat, sollen ECB und FRED trotzdem seeden.
\end{definition}

Das Skript \code{query\_db.py} setzt das um, indem jeder
Vendor-Seeder in einem separaten \code{session\_scope} läuft. Ein
Fehler in einem Vendor wird geloggt, der nächste Vendor läuft trotzdem.

\section{Was als Nächstes?}

Die ETL-Pipeline produziert Daten in \code{prices\_daily} und
\code{macro\_observations}. Im nächsten Kapitel sehen wir, wie diese
Daten in einem \emph{Jupyter-Notebook} visualisiert werden können.
